% Part: incompleteness
% Chapter: theories-computability
% Section: q-is-ce

\documentclass[../../../include/open-logic-section]{subfiles}

\begin{document}

\olfileid[fr]{inc}{tcp}{qce}

\olsection{$\Th{Q}$ est complète parmi les ensembles effectivement énumérables}

\begin{thm}
L'ensemble des codes des théorèmes de~$\Th{Q}$ est
!!{c.e.}, mais non décidable. Il est même complet,
pour la réduction plusieurs-vers-un, parmi les
ensembles !!{c.e.}s.
\end{thm}

\begin{proof}
L'ensemble des codes des théorèmes de~$\Th{Q}$ est
!!{c.e.} : ce sont les codes~$y$ pour lesquels il
existe un témoin de preuve~$x$ dans~$\Th{Q}$,
autrement dit, en identifiant la théorie à cet
ensemble de codes,
\[
Q = \Setabs{y}{\lexists[x][\Prf[\Th{Q}](x,y)]}.
\]
La relation $\Prf[\Th{Q}](x,y)$ est calculable,
même récursive primitive; la projection existentielle
d'une telle relation est effectivement énumérable.
La non-décidabilité a été démontrée à la
\olref[req][und]{sec}.

Être complet parmi les ensembles !!{c.e.}s revient
ici à montrer $K\leq_m Q$, où
$K=\Setabs{x}{!A_x(x)\downarrow}$ est le problème
diagonal de l'arrêt. Le prédicat de Kleene $T(e,x,s)$
étant récursif primitif, une !!{formula}~$!A_T$ le
représente dans~$\Th{Q}$. Pour tout~$x$, nous avons
\begin{align*}
x \in K & \liff \lexists[s][T(x,x,s)] \\
& \liff \lexists[s][(\Th{Q} \Proves !A_T(\num x,\num x, \num s))]
  \\
& \liff \Th{Q} \Proves \lexists[s][!A_T(\num x, \num x, s)].
\end{align*}
Expliquons le dernier sens réciproque. Si
$\Th{Q}\Proves\lexists[s][!A_T(\num x,\num x,s)]$,
la correction de~$\Th{Q}$ dans le modèle standard
donne un entier naturel~$n$ tel que
$\Sat{N}{!A_T(\num x,\num x,\num n)}$.
Si $T(x,x,n)$ était fausse, la condition négative
de représentation de~$T$ donnerait
$\Th{Q}\Proves\lnot !A_T(\num x,\num x,\num n)$.
La correction de~$\Th{Q}$ rendrait alors vraie
dans~$\Struct{N}$ la négation de cette instance,
contradiction. Donc $T(x,x,n)$ est vraie et
$x\in K$.

Enfin, à partir de~$x$, on peut calculer le code
de l'!!{sentence}
$\lexists[s][!A_T(\num x,\num x,s)]$.
La fonction qui envoie~$x$ sur ce code est une
réduction plusieurs-vers-un de~$K$ à l'ensemble
des codes des théorèmes de~$\Th{Q}$.
\end{proof}

\begin{explain}
La réduction associe à chaque entrée un énoncé du langage de
l'arithmétique. Le prédicat de Kleene intervient dans le
raisonnement sur la calculabilité. C'est sa formule
représentante $!A_T$ qui figure dans l'énoncé dont on code
la démonstrabilité.
\end{explain}

\end{document}
